\section{Exterior estimates and exhaustion by profiles}
\label{sec:profiles}

The fast projection identifies the possible local limits. It does not by
itself exclude norm at an internal scale that tends to zero in physical
coordinates while tending to infinity in the coordinates
\(y=h^{1/3}(b-c)\). An exterior estimate excludes precisely this possibility.
We prove that estimate simultaneously with the kernel theorem, by induction
on the size of the color algebra.

\begin{proposition}[Exterior estimate]\label{prop:exterior-estimate}
For each \(n\ge2\) there are constants \(c_n>0\) and \(R_n<\infty\) such that
\begin{equation}\label{eq:exterior-estimate}
 \norm{rD_u(h,0)F}\ge c_n\norm{F},
 \qquad
 \supp F\subset\{h^{1/3}r>R_n\},
 \qquad r=|x|,
\end{equation}
for every \(h>0\) and every compactly supported gauge-invariant graph-domain
vector for the \(\SU(n)\) operator. No rotation invariance is required.
\end{proposition}

The proof of Proposition~\ref{prop:exterior-estimate} appears below. At size
\(N\), until that proof is complete, we assume only its conclusions and
\(\dim\cK_n=1\) for \(2\le n<N\). We set \(\cK_1=\C\). Every block in a
proper separated-center chart has size less than \(N\), so all the local
arguments preceding the proof use only these induction hypotheses. At
\(N=2\) the proper blocks are scalar, and there is no nontrivial hypothesis.
In particular, we do not assume that \(\cK_N\) is nonzero or finite
dimensional.

\subsection{Intermediate scales}

\begin{lemma}[Quantitative exclusion of intermediate scales]
\label{lem:intermediate-scales}
Suppose \eqref{eq:exterior-estimate} holds at size \(n\). Let \(F_h\) have the
gauge symmetry of that block, with
\(\norm{F_h}\le B\) and \(\norm{D_u(h,0)F_h}\le K\). Both norms may include
arbitrary spectator variables or a spectator Hilbert space. There is an
absolute cutoff constant \(C\) such that, whenever \(L\ge1\) and
\(M>e^LR_n\),
\begin{equation}\label{prof:intermediate-bound}
 \left(\int_{Mh^{-1/3}<r<\varepsilon}|F_h|^2\right)^{1/2}
 \le c_n^{-1}\left(e^L\varepsilon K+\frac{CB}{L}\right).
\end{equation}
In particular,
\begin{equation}\label{eq:intermediate-scales}
 \lim_{\substack{\varepsilon\downarrow0\\M\to\infty}}
 \limsup_{h\to\infty}
 \int_{Mh^{-1/3}<r<\varepsilon}|F_h|^2=0.
\end{equation}
The outer limit is directed: after choosing an error tolerance, one may
choose \(M_0\) and \(\varepsilon_0\) so that the bound holds for every
\(M\ge M_0\) and \(0<\varepsilon\le\varepsilon_0\).
\end{lemma}

\begin{proof}
The assertion is immediate if the indicated interval is empty. Otherwise
choose \(0\le\chi\le1\), equal to one on that interval, supported in
\[
 Me^{-L}h^{-1/3}<r<e^L\varepsilon,
 \qquad |r\chi'(r)|\le C/L.
\]
One obtains such a function by multiplying two fixed smooth cutoffs in
\(\log r\), each with transition length \(L\). Its support satisfies the
hypothesis of \eqref{eq:exterior-estimate}. The Clifford normalization gives
\(\norm{[D_u,\chi](x)}_{\mathrm{op}}=|\nabla\chi(x)|/\sqrt2\). Thus
\[
 c_n\norm{\chi F_h}
 \le\norm{r\chi D_u(h,0)F_h}+
       \norm{r[D_u,\chi]F_h}
 \le e^L\varepsilon K+CB/L.
\]
Approximation in the graph norm extends this computation beyond smooth
vectors. The same estimate with spectators follows first for finite sums
of spectator vectors, by squaring the exterior inequality and integrating,
and then by approximation.

For a given \(\delta>0\), first choose \(L\) so that
\(CB/(c_nL)<\delta/2\), then choose \(M_0>e^LR_n\), and finally choose
\(\varepsilon_0\) so that \(e^L\varepsilon_0K/c_n<\delta/2\). This proves
the stated order and uniformity of the limits.
\end{proof}

\subsection{Norm measures and local compactness}

Let \(\pi\) denote the quotient by the compact gauge group. The squared
norm measure of a vector \(\psi\) is the pushforward under \(\pi\) of
\(|\psi(x)|^2\,\dd x\). Bounded sequences of these finite measures have
subsequences converging against continuous compactly supported functions.
The quotient is a locally compact separable metric space: for instance,
the distance between two orbits is the minimum Euclidean distance between
their representatives. Closed bounded quotient sets are compact.

For an unordered partition \(C=(n_1,\ldots,n_k)\) with \(k\ge2\), let
\(\mathcal P_C\) be the profile space of Definition~\ref{def:profile-space}.
Its labeled center space is
\[
 \mathcal Z_C=\left\{(c_1,\ldots,c_k)\in(\R^9)^k:
                   \sum_\alpha n_\alpha c_\alpha=0\right\},
 \qquad |c|^2=\sum_\alpha n_\alpha|c_\alpha|^2,
\]
of dimension \(d_C=9(k-1)\). We use the flat measure induced by this metric,
with the fixed normalization for permutations of equal-sized blocks in
Definition~\ref{def:profile-space}. In particular, the quotient norm below
has no extra orbit-volume factor.

\begin{proposition}[Profiles on proper strata]
\label{prop:profile-compactness}
Assume the induction hypotheses at all sizes less than \(N\). Fix \(m\ge0\),
let \(h_j\to\infty\), and suppose that \(\psi_j\) is a bounded sequence of
gauge-invariant vectors with locally bounded \(D_u(h_j,m)\) graph norms.
Here local boundedness means that
\(\norm{D_u(h_j,m)(a\psi_j)}\) is bounded for every smooth compactly
supported gauge-invariant multiplier \(a\).
After a subsequence there is a profile \(F_C\in\mathcal P_C\) on each proper
stratum such that the limiting norm measure restricted to that stratum is
exactly its squared norm density. Moreover, for every test
\(f\in\mathcal T_C\) of Proposition~\ref{prop:test-lifts},
\begin{equation}\label{prof:test-overlap}
 \ip{\psi_j}{I_{h_j,u,m}f}\longrightarrow\ip{F_C}{f}.
\end{equation}
The extraction can be made simultaneously for any finite list of
sequences, and then the analogous restrictions of all their polarized
norm measures and test overlaps hold. Parity and rotation symmetries of
the original sequences pass to the profiles under the identifications of
Lemma~\ref{lem:chart-compatibility}.
\end{proposition}

\begin{proof}
Local compactness will construct a kernel-valued profile. The
intermediate-scale estimate is then needed to show that it accounts for
all norm assigned to its stratum. We carry out these two steps before
gluing the local profiles and treating inner products.

\smallskip\noindent\emph{Construction of a local profile.}
On a compact set the fixed mass term is a bounded multiplier, uniformly
in \(h_j\).
The raw estimates of Lemma~\ref{lem:raw-estimates} therefore imply that any
limiting norm measure is supported on commuting configurations. Indeed,
on a compact subset where
\(\sum_{i<j}|[x_i,x_j]|^2\ge a>0\), its mass is \(O(h_j^{-1}/a)\).

Fix a smaller compact set in a separated-center chart, and first multiply
the states by a smooth cutoff supported in the chart and equal to one on
a neighborhood of that set. Write the resulting states again as
\(\psi_j\). All the estimates that follow are uniform on this fixed set.
Lemma~\ref{lem:fast-projection} gives
\[
 \psi_j-P_u\psi_j\longrightarrow0,
 \qquad P_u\psi_j=\chi_u(b)\eta_j(b),
 \qquad \norm{D_u^s(h_j,m)P_u\psi_j}\le C.
\]
The normalized fast section is even and has bounded derivatives as a
vector in the fast Hilbert space. The slow Clifford multipliers act on
the other factor. Consequently the Leibniz identity reads
\begin{equation}\label{prof:section-removal}
 D_u^s\bigl(\chi_u\eta_j\bigr)
 =\chi_uD_u^s\eta_j+
       \sum_\ell a_\ell(\partial_{b_\ell}\chi_u)\eta_j,
\end{equation}
where the \(a_\ell\) are constant slow Clifford matrices. The last term
has norm at most \(C\norm{\eta_j}\). In particular, no factor of \(h_j\)
is introduced when removing the fast section: the large internal bracket
multiplier acts only on slow fermions. Subtracting the bounded mass
multiplier gives
\(\norm{D_u^s(h_j,0)\eta_j}\le C\).
Covariance of \(\chi_u\) means that \(\eta_j\) has the separate gauge
symmetry in each internal \(\SU(n_\alpha)\) factor.

Put \(w_\alpha=b_\alpha-c_\alpha\Id_{n_\alpha}\). On the slow tensor
product, with all graded signs included,
\[
 D_u^s(h_j,0)=D_u^C(0)+\sum_{\alpha:n_\alpha>1}
                             D_{u,\alpha}(h_j,0).
\]
These summands involve disjoint variables and odd Clifford factors, so
their pairwise anticommutators vanish. For compactly supported smooth
sections integration by parts gives
\begin{equation}\label{prof:block-graph-identity}
 \norm{D_u^s(h_j,0)\eta_j}^2
 =\norm{D_u^C(0)\eta_j}^2+
       \sum_{\alpha:n_\alpha>1}
                         \norm{D_{u,\alpha}(h_j,0)\eta_j}^2.
\end{equation}
The identity extends by closure. One can see this directly here by
approximating after the compact slow cutoff; every potential is bounded
on its support for fixed \(h_j\), and ordinary local smoothing applies.
In particular, each internal graph norm is bounded and
\(\norm{\nabla_c\eta_j}\le C\), since
\(\norm{D_u^C(0)g}^2=\tfrac12\norm{\nabla_cg}^2\).

Rescale the internal variables unitarily by \(y_\alpha=h_j^{1/3}w_\alpha\),
and denote the resulting function by \(\widetilde\eta_j(c,y)\). Then
\begin{equation}\label{prof:rescaled-internal-bounds}
 \norm{D_{u,\alpha}(1,0)\widetilde\eta_j}\le Ch_j^{-1/3},
 \qquad \norm{\nabla_c\widetilde\eta_j}\le C.
\end{equation}
For an internal cutoff \(\zeta\) of fixed compact support, write the
constant-coefficient kinetic Dirac operator as
\(D_{u,\alpha}(1,0)\) minus its bracket multiplier. That multiplier is
bounded on \(\supp\zeta\). The kinetic norm identity and the bounded
commutator with \(\zeta\) give
\[
 \norm{\nabla_{y_\alpha}(\zeta\widetilde\eta_j)}
 \le C_\zeta\bigl(\norm{D_{u,\alpha}(1,0)\widetilde\eta_j}
                          +\norm{\widetilde\eta_j}\bigr).
\]
All spectators are integrated in these norms. Combining these inequalities
for the finitely many blocks with the center derivative estimate gives a
bounded ordinary \(H^1\) norm on every compact subset of the joint
\((c,y)\) space. Its Clifford fiber is finite dimensional. The usual
Rellich theorem therefore gives, by a diagonal subsequence, strong local
\(L^2\) convergence to a function \(\eta_\infty(c,y)\). Compactness has
been applied jointly in all these variables, rather than to functions
with an uncontrolled infinite-dimensional spectator fiber.

Fatou's lemma gives the global \(L^2\) bound for this local chart limit.
Testing \eqref{prof:rescaled-internal-bounds} against compactly supported
smooth functions shows distributionally that
\(D_{u,\alpha}(1,0)\eta_\infty=0\) for each \(\alpha\). Fubini's theorem
and a countable dense family of internal tests show that almost every
slice in the remaining variables is in the corresponding distributional
\(L^2\) kernel. Lemma~\ref{lem:domains} identifies this with
\(\cK_{n_\alpha}\). Equivalently, the orthogonal projection onto each
internal kernel fixes \(\eta_\infty\); the product of these commuting
projections identifies its coefficient space as
\(\bigotimes_\alpha\cK_{n_\alpha}\). This constructs the local profile.

\smallskip\noindent\emph{Exact norm on the stratum.}
We now identify how much limiting norm is assigned to the stratum
\(w=0\). The replacement of \(\psi_j\) by \(P_u\psi_j\) changes integrals against
bounded multipliers by \(o(1)\). The physical fast coordinates tend to zero
in norm measure, since \(z=t/\sqrt{h_j}\) and
\(\norm{t\psi_j}\le C\). On bounded \(y\) sets, smoothness also gives
\(\chi_u(c+h_j^{-1/3}y)-\chi_u(c)\to0\) uniformly in fast Hilbert norm.
It remains to decide whether any internal norm escapes all bounded
\(y\) sets while still reaching \(w=0\).

Here is a precise measure comparison. Take a nonnegative continuous center
test \(a(c)\) of compact support in a smaller chart. Let
\(0\le\kappa_\varepsilon(w)\le1\) be a product of radial internal cutoffs,
equal to one when every \(|w_\alpha|\le\varepsilon\) and zero when some
\(|w_\alpha|\ge2\varepsilon\). Choose them decreasing as
\(\varepsilon\downarrow0\). In testing the quotient measure, include a fixed
physical fast cutoff equal to one near \(z=0\), so that the product extends
continuously by zero inside the outer chart. Its removal changes the limit
by zero, because \(\norm{t\psi_j}\le C\) and \(z=t/\sqrt{h_j}\).
The slice density has already been absorbed unitarily, and
\(\int|\chi_u(b,t)|^2\,\dd t=1\); thus, after removal of the fast cutoff,
these tests of the projected states are exactly the corresponding slow
integrals. For each fixed \(M\), strong local convergence
identifies the integrals on \(\{|y_\alpha|\le M\text{ for all }\alpha\}\).
The difference between such an integral and the integral with
\(\kappa_\varepsilon\) is bounded above, apart from this local convergence
error, by
\begin{equation}\label{prof:measure-tail-bound}
 \norm{a}_\infty\sum_{\alpha:n_\alpha>1}
 \int_{Mh_j^{-1/3}<|w_\alpha|<2\varepsilon}|\eta_j|^2.
\end{equation}
Each summand satisfies Lemma~\ref{lem:intermediate-scales}, with all other
variables as spectators, because of \eqref{prof:block-graph-identity}.
For the converse comparison, the bounded \(y\) set lies where
\(\kappa_\varepsilon=1\) for all sufficiently large \(j\). Thus its limiting
integral is a lower bound for the one with \(\kappa_\varepsilon\).

To record the order of limits explicitly, put
\[
 J_\varepsilon=\int a(c)\kappa_\varepsilon(w)\,\dd\mu,
 \qquad
 A_M=\int_{\max_\alpha|y_\alpha|\le M}
                         a(c)|\eta_\infty(c,y)|^2\,\dd c\,\dd y.
\]
For every fixed \(L\ge1\) and
\(M>e^L\max_\alpha R_{n_\alpha}\), the preceding comparisons and
\eqref{prof:intermediate-bound}, after \(j\to\infty\), give
\[
 A_M\le J_\varepsilon\le A_M+
 \norm{a}_\infty\sum_{\alpha:n_\alpha>1}
 c_{n_\alpha}^{-2}\left(2e^L\varepsilon K_\alpha+CB/L\right)^2,
\]
where \(K_\alpha\) are the block graph bounds in
\eqref{prof:block-graph-identity}. Decreasing \(\varepsilon\) restricts
\(J_\varepsilon\) to \(w=0\), by continuity from above and concentration
of the fast coordinates at zero. Next let \(M\to\infty\), keeping its
lower threshold, and finally let \(L\to\infty\). The error vanishes and
\(A_M\) exhausts the profile norm. Thus the restricted measure is exactly
\(|F_C(c)|^2\,\dd c\). If every block is scalar, the sums and internal
cutoffs are absent and strong local center convergence gives the same
conclusion directly.
This argument also covers simultaneous fragmentation in several blocks,
since their error sets are bounded by the finite sum in
\eqref{prof:measure-tail-bound}.

\smallskip\noindent\emph{Compatibility and inner products.}
Local strong convergence and approximation by bounded \(y,t\) sets,
followed by Proposition~\ref{prop:test-lifts}, give the local version of
\eqref{prof:test-overlap}. It also identifies the profile independently
of the outer chart cutoff. If two cutoffs \(\lambda,\widetilde\lambda\)
equal one near the same smaller center support, then
\(\| (\lambda-\widetilde\lambda)I_{h_j,u,m}f\|\to0\) for every test
supported there: its leading lift concentrates in that neighborhood,
and its correction tends to zero in norm. Pairing with the original
bounded sequence proves that the two localizations have identical test
overlaps. These identify the same profile coefficient, including its phase.

There are finitely many partition types and a countable exhaustion of each
stratum by relatively compact charts. Repeating the extraction diagonally
gives profiles on all of them. On overlapping charts,
Lemma~\ref{lem:chart-compatibility} identifies the leading tests and hence,
by their overlaps, the local profiles. These consequently descend with
the required permutation action. Exhaustion of the stratum gives the
global norm bound, and the local overlaps give
\eqref{prof:test-overlap} for every compactly supported profile test.

For a finite list \(\psi_j^1,\ldots,\psi_j^q\), carry out the same
extractions for the list and for the finitely many sums and imaginary
sums needed for complex polarization. Equivalently, take weak limits of
the matrix of measures with entries
\(\pi_*(\ip{\psi_j^a(x)}{\psi_j^b(x)}\,\dd x)\). Positivity of this matrix
and Cauchy--Schwarz bound every entry by its diagonal measures. Linearity
of the local limits then gives, on each stratum, the entry
\(\ip{F_C^a(c)}{F_C^b(c)}\,\dd c\). This proves the polarized assertion.
The same leading identifications commute with the indicated symmetries,
which proves the final assertion.
\end{proof}

\subsection{Passing the charge to the center space}

\begin{lemma}[Removal of collisions]\label{lem:collision-removal}
Let \(S_C\subset \mathcal Z_C\) be the union of the collision subspaces
\(c_\alpha=c_\beta\). If \(F,G\in L^2(\mathcal Z_C)\), with the profile coefficient
fiber, satisfy \(D_u^C(m)F=G\) distributionally on \(\mathcal Z_C\setminus S_C\),
then the same identity holds on \(\mathcal Z_C\). In particular, \(F\) is in the
full Euclidean graph domain. This remains true on the subspace with the
permutation symmetry of the profile space.
\end{lemma}

\begin{proof}
Each collision subspace has codimension nine, also after imposing the
trace-zero constraint. Let \(\rho_\varepsilon\) be the product of smooth
cutoffs that vanish at distance at most \(\varepsilon\) from these finitely
many subspaces and equal one at distance at least \(2\varepsilon\). On
each compact set \(K\subset \mathcal Z_C\), tubular coordinates give
\begin{equation}\label{prof:collision-capacity}
 \norm{\nabla\rho_\varepsilon}_{L^2(K)}^2
 \le C_K\varepsilon^{-2}\varepsilon^9=C_K\varepsilon^7.
\end{equation}
Intersections of collision subspaces cause no difficulty, since the
derivative of the finite product is bounded by the sum of the individual
derivatives.

For a smooth compactly supported test \(\varphi\), apply the given
distributional identity to \(\rho_\varepsilon\varphi\). The commutator
error has absolute value at most
\[
 2^{-1/2}\norm{F}_{L^2(K)}\norm{\varphi}_\infty
                    \norm{\nabla\rho_\varepsilon}_{L^2(K)},
\]
which tends to zero. All other terms converge by dominated convergence;
the mass multiplier is bounded on \(K\). This proves the global
distributional equation. The maximal and minimal graph domains coincide
by the cutoff and local ellipticity argument of Lemma~\ref{lem:domains},
which applies verbatim to this abelian operator. Tests with values in a
finite-dimensional coefficient subspace suffice, followed by density if
necessary. For permutation-equivariant vectors, averaging the tests over
the finite permutation group gives the same conclusion. There is thus no
additional collision boundary condition.
\end{proof}

\begin{proposition}[Liminf for the center charges]
\label{prop:profile-liminf}
Under the hypotheses of Proposition~\ref{prop:profile-compactness}, suppose
in addition that \(\norm{D_u(h_j,m)\psi_j}\) is globally bounded. On a
subsequence realizing its lower limit, the profiles satisfy
\begin{equation}\label{eq:profile-liminf}
 \sum_{C\text{ proper}}\norm{D_u^C(m)F_C}^2
       \le\liminf_{j\to\infty}\norm{D_u(h_j,m)\psi_j}^2.
\end{equation}
The derivatives on the left are the full-space graph derivatives.
\end{proposition}

\begin{proof}
Choose any finite family of profile tests \(g_C\in\mathcal T_C\) and put
\(T_jg=\sum_C I_{h_j,u,m}g_C\). Proposition~\ref{prop:test-lifts} and
Lemma~\ref{lem:chart-compatibility} give
\[
 \norm{T_jg}^2\longrightarrow\sum_C\norm{g_C}^2,
 \qquad
 \ip{D_u(h_j,m)\psi_j}{T_jg}
 \longrightarrow\sum_C\ip{F_C}{D_u^C(m)g_C}.
\]
For the second identity, move the self-adjoint charge to the tests; its
graph-output error pairs with the bounded norm of \(\psi_j\) and tends
to zero. Cauchy--Schwarz therefore bounds the last functional by
\[
 \left(\liminf_j\norm{D_u(h_j,m)\psi_j}\right)
                 \left(\sum_C\norm{g_C}^2\right)^{1/2}.
\]
The tests are dense in the Hilbert direct sum of the profile spaces.
Riesz representation identifies this bounded functional with the direct
sum of the distributional derivatives, and proves
\eqref{eq:profile-liminf} off the collisions. Lemma~\ref{lem:collision-removal}
extends each derivative across them. This argument uses convergence for
each fixed finite test family; it requires no uniform estimate on the
whole unit ball of test functions.
\end{proof}

\pagebreak
\subsection{Hardy inequality and the current exterior estimate}

We can now prove Proposition~\ref{prop:exterior-estimate} at size \(N\),
without having used any information about \(\cK_N\).

\begin{proof}[Proof of Proposition~\ref{prop:exterior-estimate} at size \(N\)]
Fix a compact annulus \(a\le r\le b\), where \(0<a<b<\infty\). We first
claim that every sequence of unit vectors supported there, with
\(h_j\to\infty\), satisfies
\begin{equation}\label{prof:fixed-annulus-bound}
 \liminf_j\norm{rD_u(h_j,0)F_j}\ge\frac{5}{2\sqrt2}.
\end{equation}
It suffices to treat a subsequence for which the norms on the left are
bounded. Because \(r\ge a\), its unweighted graph norms are bounded.
Every limiting norm measure is supported on commuting configurations in
this annulus. A traceless commuting configuration with only one joint
center is zero. Thus only proper strata occur here, and
Proposition~\ref{prop:profile-compactness}, together with additivity of the
measure over the finitely many strata, gives
\begin{equation}\label{prof:annulus-norm-exhaustion}
                  \sum_{C\text{ proper}}\norm{F_C}^2=1.
\end{equation}

The profiles of \(rF_j\) are \(|c|F_C\). This follows either from the
local rescalings or from \eqref{prof:test-overlap}, since \(r\) is bounded
on the common support and restricts to \(|c|\) on the stratum. The
commutator with \(r\) is bounded, so Proposition~\ref{prop:profile-liminf}
applies to this sequence as well.

For a reference for the classical Euclidean Hardy inequality, see
\cite[equation~(1.1), with $p=2$]{FrankSeiringer2008}.
For completeness, completing a square in Euclidean dimension \(d>2\)
gives, for smooth compactly supported \(g\) away from zero,
\[
 0\le\left\|\nabla g+\frac{d-2}{2}\frac{c}{|c|^2}g\right\|^2
 =\norm{\nabla g}^2-\frac{(d-2)^2}{4}\norm{g/|c|}^2.
\]
Cutoffs and closure extend this Hardy inequality to \(H^1\). Applied on
\(\mathcal Z_C\), with its orthonormal trace coordinates and \(d_C=9(k-1)\ge9\),
it yields
\[
 \norm{D_u^C(0)(|c|F_C)}
 =2^{-1/2}\norm{\nabla_c(|c|F_C)}
 \ge\frac{d_C-2}{2\sqrt2}\norm{F_C}
 \ge\frac7{2\sqrt2}\norm{F_C}.
\]
The full-space graph statement needed here is supplied by
Lemma~\ref{lem:collision-removal}. Hence
\eqref{eq:profile-liminf} and \eqref{prof:annulus-norm-exhaustion} imply
\(\liminf_j\norm{D_u(h_j,0)(rF_j)}\ge7/(2\sqrt2)\).
Finally, \(\norm{[D_u,r]F_j}=1/\sqrt2\), proving
\eqref{prof:fixed-annulus-bound}.

Fix a logarithmic width \(A>0\). Since \(5/(2\sqrt2)>1\), a contradiction
sequence in the fixed annulus \(1<r<e^A\) shows that there is \(H_A<\infty\)
such that its annular inequality \(\norm{rD_u(h,0)F}\ge\norm{F}\) holds
for every \(h\ge H_A\). Dilation \(x=s\xi\) replaces the coupling by
\(hs^3\) and leaves the weighted operator \(rD_u\) unchanged. Therefore
\begin{equation}\label{prof:scaled-annulus-bound}
 \norm{rD_u(h,0)F}\ge\norm{F}
 \quad\text{if}\quad
 \supp F\subset\{s<r<e^As\},\quad h^{1/3}s\ge H_A^{1/3}.
\end{equation}

We use the IMS localization principle~\cite[Lemma~3.1]{Simon1983},
here in its weighted first-order form. Choose smooth functions
\(\chi_j(\log r)\) forming a squared partition of
unity, with supports in intervals of length \(A\), bounded overlap, and
\(\sum_j|\chi_j'|^2\le C/A^2\). For example, normalize translates of a
fixed bump by the square root of the sum of their squares, then rescale
the real variable. The Clifford symbol and
\(\sum_j\chi_j\nabla\chi_j=0\) give the exact localization identity
\begin{equation}\label{prof:weighted-ims}
 \sum_j\norm{rD_u(\chi_jF)}^2
 =\norm{rD_uF}^2+
     \frac12\int r^2\sum_j|\nabla\chi_j|^2|F|^2
 \le\norm{rD_uF}^2+\frac{C}{A^2}\norm{F}^2.
\end{equation}
Fix \(A\) large enough that \(C/A^2\le1/2\), and only then set
\(R_N>e^AH_A^{1/3}\). Every member of the partition meeting
\(\supp F\subset\{h^{1/3}r>R_N\}\) is supported in an annulus satisfying
\eqref{prof:scaled-annulus-bound}: its lower radius is at least \(e^{-A}\)
times a radius where it meets \(F\). Summing that bound and using
\eqref{prof:weighted-ims} proves
\(\norm{rD_uF}^2\ge\tfrac12\norm{F}^2\). Taking, for example,
\(c_N=1/\sqrt2\) proves the assertion. The threshold can depend on \(N\).
All identities initially use compact smooth vectors and extend to the
stated domain by Lemma~\ref{lem:domains}.
\end{proof}

\subsection{The full-collapse profile}

\begin{corollary}[Exhaustion for tight sequences]
\label{prof:full-exhaustion}
Under the hypotheses of Proposition~\ref{prop:profile-compactness}, after
Proposition~\ref{prop:exterior-estimate} has been proved at size \(N\),
there is also a full-collapse profile \(F_0\in\cK_N\). Its squared norm is
the mass of the limiting norm measure at zero. If the squared norm
measures are tight, then
\begin{equation}\label{prof:full-norm-exhaustion}
 \lim_j\norm{\psi_j}^2
       =\norm{F_0}^2+\sum_{C\text{ proper}}\norm{F_C}^2.
\end{equation}
For any finite simultaneously extracted list, one likewise has
\begin{equation}\label{prof:polarized-exhaustion}
 \lim_j\ip{\psi_j^a}{\psi_j^b}
 =\ip{F_0^a}{F_0^b}+
       \sum_{C\text{ proper}}\ip{F_C^a}{F_C^b}.
\end{equation}
The space \(\cK_N\) is allowed to have unknown dimension in this statement.
\end{corollary}

\begin{proof}
Multiply by a fixed compact radial cutoff equal to one near zero.
The resulting undeformed graph norms are bounded, since the mass term is
bounded there. Unitary whole-configuration rescaling by \(h_j^{1/3}\)
gives bounded \(L^2\) functions whose \(D_u(1,0)\) outputs tend to zero.
Local ellipticity and Rellich compactness give a strong local limit.
Fatou's lemma makes that limit globally square integrable, and
Lemma~\ref{lem:domains} puts it in \(\cK_N\). Lemma~\ref{lem:intermediate-scales},
now available at size \(N\), identifies its norm with exactly the atom at
zero by the same inner-ball and shrinking-physical-ball comparison used
in \eqref{prof:measure-tail-bound}. This construction is independent of
the chosen cutoff near zero. It extracts one vector in \(\cK_N\); it
does not require compactness of its unit ball or prior finite
dimensionality.

The commuting quotient is the disjoint union of zero and the proper
strata. Tightness preserves total mass in the weak limit: approximate
the constant function one by compact cutoffs and use the uniform tail
bound. Adding the restricted measures gives
\eqref{prof:full-norm-exhaustion}. Applying this to the finite linear
combinations used for polarization gives
\eqref{prof:polarized-exhaustion}. Rotation and parity pass to \(F_0\)
directly under radial cutoffs and dilation.
\end{proof}
